% \vspace{-2ex}
\section{Related Work}
% \vspace{-1ex}
\paragraph{Quantization-Aware Training.}
In recent years, model quantization has been a topic of great interest in the deep learning community due to neural networks continuously scaling exponentially in terms of compute. Neural network quantization approaches can be broadly categorized into: \acrfull{PTQ} and \acrfull{QAT}. Though \acrshort{PTQ} \cite{dfq, bannerposttraining, nagel2021white, adaround} is faster and does not rely on whole training data, it yields significant performance degradation at low-bit precision. \acrshort{QAT} is the focus of our work, and it has been well-studied in literature \cite{Gong2019DSQ, pact2018, zhou2017incremental, lsq, liu2022nonuniform, EWGS}.

The \acrshort{STE} is a de-facto method for backpropagating through the non-differentiable rounding function in \acrshort{QAT}. The effectiveness of \acrshort{STE} has been the point of argument in recent literature. Ajanthan \etal \cite{pmlr-v130-ajanthan21a} proposed a mirror descent formulation for neural network quantization and established the connection between \acrshort{STE} approximation and mirror descent framework for constrained optimization. Lee \etal~\cite{EWGS} showed that the \acrshort{STE} leads to bias in gradients and proposed gradient scaling by the distance of latent weights from quantization boundaries. Gong \etal \cite{Gong2019DSQ} attempt to mitigate the issues caused by \acrshort{STE} for low-bit quantization by using differentiable hyperbolic tangent functions to simulate the rounding function in the backward pass. Similar to that, Yang \etal~\cite{quantization_nets} approximate the rounding function using smooth sigmoid functions to address the gradient bias in \acrshort{STE}. 

Recent works~\cite{defossez2022differentiable, nagel2022overcoming} have identified oscillation as a side-effect of \acrshort{STE} approximation during \acrshort{QAT}. Defossez \etal~\cite{defossez2022differentiable} proposed additive Gaussian noise to mimic the quantization noise and replace it as quantization operation during \acrshort{QAT} to prevent weight oscillations and biased gradients resulting from \acrshort{STE}. Nagel \etal \cite{nagel2022overcoming} also showed that weight oscillation seriously impacts \acrshort{QAT} performance, specifically on efficient networks comprising depth-wise convolutions due to \acrshort{STE} approximation of rounding function. They propose to constrain the latent weights to avoid oscillation by either regularizing them to their quantized states or by freezing them. Recently, Liu \etal~\cite{liu2023oscillation} studied the issue of oscillations on vision transformers and proposed fixed scale factors for weight quantization and query-key re-parameterization to mitigate the negative influence of oscillation. However, their proposed approach is specifically targeted at solving oscillation phenomenon on transformers architecture. 

\paragraph{Quantization of  Object Detectors.}
The majority of existing neural network quantization literature focuses on the image classification task rather than real-world downstream tasks such as object detection or semantic segmentation. 
Some recent works do explore the quantization of object detectors to improve the efficiency of these models. Jacob \etal~\cite{jacob2018cvpr} proposed a quantization scheme using integer-only arithmetic and performing object detection on \coco{} dataset but the approach is only effective for 8-bit quantization. Li et al.~\cite{li2019fully} observed training instability during the quantized fine-tuning of RetinaNet and propose mutliple solutions specific to RetinaNet architecture. These solutions are not applicable on more efficient object detectors like \yolo{}. Ding \etal~\cite{ding2019req} proposed ADMM based weight quantization framework for \yolo{3} but do not address the issue of oscillations and activation quantization.  Due to the discrete nature of quantization functions, gradient estimates are known to be noisy, which affects \acrfull{SGD} updates. To overcome that, Zhuang \etal~\cite{zhuang2020training} proposed to utilize a full precision auxiliary module to enable stable training of quantized object detectors. Furthermore, Zhuang \etal~\cite{chen2021aqd} proposed multi-level \acrfull{BN} to accurately calculate batch statistics for each pyramid level in RetinaNet~\cite{lin2017focal} and FCOS~\cite{tian2019fcos}. The proposed approach is specific to pyramid-level architecture in RetinaNet and FCOS that share \acrshort{BN} layers at different levels of the pyramid, but not applicable on more recent efficient SOTA \yolo{} 
models. In this paper, we identify the gap in the recent literature on quantized object detection and introduce state-of-the-art quantized object detectors using \yoloa{}~\cite{yolov5} and \yolob{}\cite{wang2022yolov7} variants. 

